\chapter{Introduction}

%\label{chap:GettingStarted}
\section{Overview}
With the increase of people using credit cards in their daily lives, credit card companies should 
take special care of the security and safety of the customers. According to (Credit card statistics 
2021), the number of people using credit cards worldwide was 2.8 billion in 2019; also, 70% of 
those users own a single card.
Reports of Credit card fraud in the U.S. rose by 44.7% from 271,927 in 2019 to 393,207 words 
in 2020. There are two kinds of credit card fraud, and the first is having a credit card account opened 
under your name by an identity thief. Reports of this fraudulent behaviour increased 48% from 2019 
to 2020. The second type is when an identity thief uses an existing account you created, usually by 
stealing the information on the credit card. Reports on this type of Fraud increased 9% from 2019 
to 2020 (Daly, 2021). Those statistics caught  We's attention as the numbers have  increased 
drastically and rapidly throughout the years, which motivated We to resolve the issue analytically 
by using different machine learning methods to detect fraudulent credit card transactions within 
numerous transactions.








\section{ Project goals}
The main aim of this project is the detection of fraudulent credit card transactions, as it is 
essential to figure out the fraudulent transactions so that customers do not get charged for the 
purchase of products that they did not buy. Fraudulent Credit card transactions will be detected 
with multiple ML techniques. Then, a comparison will be made between the outcomes and results 
of each method to find the best and most suited model for detecting fraudulent credit card 
transactions; graphs and numbers will also be provided. In addition, it explores previous literature 
and different techniques used to distinguish Fraud within a dataset.

\textbf{Research question: } What machine learning model is most suited for detecting fraudulent credit 
card transactions?



